%!TEX root = ./syllabus-cis4930-cis5930.tex


\subsection*{Course Calendar}

\begingroup
\small
\renewcommand{\arraystretch}{1.2}
\setlength{\tabcolsep}{4pt}
\begin{longtable}{|>{\centering\arraybackslash}p{0.08\textwidth}|
                >{\centering\arraybackslash}p{0.10\textwidth}|
                p{0.56\textwidth} |
                >{\raggedright\arraybackslash}p{0.18\textwidth}|}
\hline
\textbf{Class} & \textbf{Date} & \textbf{Topic} & \textbf{Notes} \\
\hline
\endhead
\multicolumn{4}{|c|}{\parbox{0.95\textwidth}{\centering\textbf{Introduction}}} \\
\hline
1
  & 8/24\newline (M)
  & Course Syllabus\newline
    Preparation for Learning Edge Networks\newline
    (XG History and Course Tools)
  & Mandatory first-day attendance.\newline \textcolor{red}{\textbf{No-shows will be dropped.}} \\
\hline
\multicolumn{4}{|c|}{\parbox{0.95\textwidth}{\centering\textbf{Don't rush to test the speed; first, let's understand the 'G' in 5G.}}} \\
\hline
2
  & 8/26\newline (W)
  & XG Physical Layer\newline
    (CDMA, OFDM, MIMO, Beamforming, mmWave)\newline
    Recommended reading: MATLAB 5G Toolbox
  & HW 1 posted on Canvas \\
\hline
3-4
  & 8/31\newline (M)\newline 9/2\newline (W)
  & XG Architecture Basics\newline
    (Generations, Mobility, RAN/Core, O-RAN, Security, Market)\newline
    Recommended reading: How to deploy ``antivirus software'' on each 5G base station?
  & 9/2 (W): HW 1 DDL at 11:59 p.m. (ET)\newline 9/2 (W): HW 2 posted on GitHub Classroom \\
\hline
-
  & 9/7\newline (M)
  & No class due to Labor Day.
  &  \\
\hline
5
  & 9/9\newline (W)
  & Integrated Sensing and Communication (ISAC)\newline
    (What it is, why it is hard, and how OFDM does it)\newline
    Recommended reading: ISAC Report
  &  \\
\hline
6
  & 9/14\newline (M)
  & Paper Practicum 1 (XG I): Inside the Base Station --- Software and Silicon\newline
    (Two student presentations, 15 min each + 5 min Q\&A)\newline
    Papers: SwapRAN; MegaStation
  &  \\
\hline
7
  & 9/16\newline (W)
  & Paper Practicum 2 (XG II): Operating the Base Station --- Efficiency, AI, Adversary\newline
    (Three student presentations, 15 min each + 5 min Q\&A)\newline
    Papers: RENC; SAGE; SNI5GECT
  & HW 2 DDL at 11:59 p.m. (ET) \\
\hline
\multicolumn{4}{|c|}{\parbox{0.95\textwidth}{\centering\textbf{You've used Wi-Fi your whole life. It's far more fun than you think --- and after today, you'll know it better than anyone.}}} \\
\hline
8-9
  & 9/21\newline (M)\newline 9/23\newline (W)
  & Wi-Fi History\newline
    (802.11b to 802.11be, PHY Waveforms, Spectrum, MIMO, OFDMA, the 60 GHz ad/ay Branch)\newline
    Recommended reading: IEEE Wi-Fi Timeline
  & 9/23 (W): HW 3 posted on GitHub Classroom \\
\hline
10
  & 9/28\newline (M)
  & Paper Practicum 3 (Wi-Fi I): Why Your Video Call Stutters on Wi-Fi\newline
    (Two student presentations, 15 min each + 5 min Q\&A)\newline
    Papers: Zhuge; AFR
  &  \\
\hline
11
  & 9/30\newline (W)
  & Paper Practicum 4 (Wi-Fi II): Inside 802.11 --- Which Knobs Can We Still Redesign?\newline
    (Three student presentations, 15 min each + 5 min Q\&A)\newline
    Papers: Law; BLADE; SMUFF
  &  \\
\hline
\multicolumn{4}{|c|}{\parbox{0.95\textwidth}{\centering\textbf{Communication Was Only the Radio's First Job. The Same Wave That Carries Your Bits Can Also See the Room.}}} \\
\hline
12
  & 10/5\newline (M)
  & Wi-Fi Sensing\newline
    (Where Is the Phone? Signal Strength, Phase, Angle with ArrayTrack, Distance with Chronos, Both at Once with SpotFi)\newline
    Recommended reading: Hands-on Wireless Sensing: CSI Feature Extraction (Tsinghua)
  &  \\
\hline
13
  & 10/7\newline (W)
  & Paper Practicum 5 (ISAC): From Reference Signals to Vital Signs\newline
    (Three student presentations, 15 min each + 5 min Q\&A)\newline
    Papers: CARTS; CommRad; 2FiA
  & HW 3 DDL at 11:59 p.m. (ET) \\
\hline
14
  & 10/12\newline (M)
  & Using 802.11 Beamforming Feedback as a Sensing Signal\newline
    (Why Beamform, How the AP Asks: NDPA and NDP, From CSI to the Compressed BFM, Angles with MUSIC, the Same Angles from the BFM)\newline
    Live demo: tracking people with a Pixel Watch 3 --- BFMScan (ACM UbiComp 2026)
  & HW 4 posted on GitHub Classroom \\
\hline
\multicolumn{4}{|c|}{\parbox{0.95\textwidth}{\centering\textbf{Don't Be Afraid of Artificial Intelligence. It's Not That Intelligent.}}} \\
\hline
15
  & 10/14\newline (W)
  & Core Models of Modern AI\newline
    (CNN, RNN, Transformer, LLM, Agent, Diffusion)\newline
    Recommended reading: Dive into Deep Learning
  &  \\
\hline
16
  & 10/19\newline (M)
  & PyTorch Tutorial: From Colab to a Trained Model\newline
    (Google Colab with a GPU, Tensors, the Training Pipeline: Data, Network, Loss, Optimizer, Train/Validate/Test; Live Demo: Choosing a Neural Network, CNN vs. RNN vs. Transformer)\newline
    Recommended reading: PyTorch: Learn the Basics
  &  \\
\hline
17
  & 10/21\newline (W)
  & Semantic Communication\newline
    (Beyond bits: task-oriented and knowledge-driven communication --- why transmit every bit when the receiver only needs the meaning?)\newline
    (One student presentation, 15 min + 5 min Q\&A)\newline
    Papers: CoOrbit
  &  \\
\hline
18
  & 10/26\newline (M)
  & Paper Practicum 6 (Semantic Communication I): The Representation Shrinks\newline
    (Three student presentations, 15 min each + 5 min Q\&A)\newline
    Papers: Empilo; MagicStream; Gemino
  &  \\
\hline
19
  & 10/28\newline (W)
  & Paper Practicum 7 (Semantic Communication II): From Pixels to Tokens, and Surviving Loss\newline
    (Three student presentations, 15 min each + 5 min Q\&A)\newline
    Papers: KDC; NIER; Artic
  & Final Project Release \\
\hline
\multicolumn{4}{|c|}{\parbox{0.95\textwidth}{\centering\textbf{The Cloud Is Far Away. The Phone Is Small. Edge Intelligence Is Deciding What Runs Where.}}} \\
\hline
20
  & 11/2\newline (M)
  & Edge AI Systems: Where Should the Model Run?\newline
    (Why On-Device AI, Memory/Compute/Energy Limits, Quantization and Mixture-of-Experts, KV Cache, Split and Collaborative Inference, Model Routing)\newline
    (One student presentation, 15 min + 5 min Q\&A)\newline
    Papers: MobileMoE
  &  \\
\hline
21
  & 11/4\newline (W)
  & Paper Practicum 8 (Edge AI I): Making AI Fit on the Device\newline
    (Three student presentations, 15 min each + 5 min Q\&A)\newline
    Papers: ElastiLM; D2MoE; KVSwap
  & HW 4 DDL at 11:59 p.m. (ET) \\
\hline
22
  & 11/9\newline (M)
  & Paper Practicum 9 (Edge AI II): Where Should Intelligence Run?\newline
    (Three student presentations, 15 min each + 5 min Q\&A)\newline
    Papers: SAIL; SMEC; Confidant
  &  \\
\hline
-
  & 11/11\newline (W)
  & No class due to Veterans Day observed.
  &  \\
\hline
\multicolumn{4}{|c|}{\parbox{0.95\textwidth}{\centering\textbf{Seeing the Room Is Not Enough. An Edge Agent Has to Act --- on Time, Safely, and Without Giving You Away.}}} \\
\hline
23
  & 11/16\newline (M)
  & Edge AI Agents: From Perception to Action\newline
    (The Agent Loop, End-to-End Latency Budgets, Choosing What to Sense, Checking Actions Before They Run, Privacy)
  &  \\
\hline
24
  & 11/18\newline (W)
  & Paper Practicum 10 (Edge Agents): From Perception to Action\newline
    (Three student presentations, 15 min each + 5 min Q\&A)\newline
    Papers: Modality Plug-and-Play; TimelyLLM; VeriSafe Agent
  &  \\
\hline
25
  & 11/23\newline (M)
  & Paper Practicum 11 (Edge Privacy): What the Signal, the Features, and the Cache Give Away\newline
    (Three student presentations, 15 min each + 5 min Q\&A)\newline
    Papers: Poison to Cure; FaceObfuscator; KV-Cloak
  &  \\
\hline
-
  & 11/25\newline (W)
  & No class due to Thanksgiving holiday.
  &  \\
\hline
26
  & 11/30\newline (M)
  & Paper Practicum 12\newline
    (Four student presentations, 15 min each + 5 min Q\&A)
  &  \\
\hline
-
  & 12/2\newline (W)
  & No class. Use the time to work on your final project.
  &  \\
\hline
\multicolumn{4}{|c|}{\parbox{0.95\textwidth}{\centering\textbf{Final Exam Week}}} \\
\hline
-
  & 12/7\newline (M)\newline 12/9\newline (W)
  & No class, no final exam.
  & 12/9 (W): Final Project DDL at 11:59 p.m. (ET) \\
\hline
\end{longtable}
\endgroup

The course calendar is tentative. Please refer to the course website for the most up-to-date lecture topics, assignment deadlines, and course notices.
